%% ============================================================
%%  econcomplex --- Technical Documentation
%%  Compile with: pdflatex econcomplex_documentation.tex
%%                pdflatex econcomplex_documentation.tex  (second pass)
%% ============================================================
\documentclass[12pt, a4paper]{article}

%% --- Packages ---------------------------------------------------
\usepackage[utf8]{inputenc}
\usepackage[T1]{fontenc}
\usepackage[english]{babel}
\usepackage{lmodern}
\usepackage{geometry}
\usepackage{hyperref}
\usepackage{xcolor}
\usepackage{booktabs}
\usepackage{longtable}
\usepackage{array}
\usepackage{amsmath}
\usepackage{amssymb}
\usepackage{listingsutf8}   % replaces listings with full UTF-8 support
\usepackage{graphicx}
\usepackage{caption}
\usepackage{subcaption}
\usepackage{float}
\usepackage{parskip}
\usepackage{enumitem}
\usepackage{fancyhdr}
\usepackage{titlesec}
\usepackage{tcolorbox}
\usepackage{mdframed}

%% --- Geometry ------------------------------------------------
\geometry{
  a4paper,
  top=2.5cm, bottom=2.5cm,
  left=2cm, right=2cm,
  headheight=14pt
}

%% --- Colors ----------------------------------------------------
\definecolor{codebg}{RGB}{245,245,245}
\definecolor{codeframe}{RGB}{200,200,200}
\definecolor{pyblue}{RGB}{31,119,180}
\definecolor{pyorange}{RGB}{214,110,0}
\definecolor{pygray}{RGB}{100,100,100}
\definecolor{tblhead}{RGB}{40,70,120}
\definecolor{noteblue}{RGB}{220,235,250}
\definecolor{noteborder}{RGB}{70,130,180}
\definecolor{warnbg}{RGB}{255,248,220}
\definecolor{warnborder}{RGB}{200,150,0}

%% --- UTF-8 mapping for lstlistings ------------------------
%% (required for accented characters inside code blocks)
\lstset{
  language=Python,
  basicstyle=\ttfamily\footnotesize,
  backgroundcolor=\color{codebg},
  frame=single,
  framerule=0.5pt,
  rulecolor=\color{codeframe},
  keywordstyle=\bfseries\color{pyblue},
  commentstyle=\itshape\color{pygray},
  stringstyle=\color{pyorange},
  numbers=left,
  numberstyle=\tiny\color{pygray},
  numbersep=8pt,
  showstringspaces=false,
  breaklines=true,
  breakatwhitespace=true,
  tabsize=4,
  captionpos=b,
  xleftmargin=1.5em,
  xrightmargin=0.5em,
  inputencoding=utf8,
  extendedchars=true,
  literate=
    {\'a}{{\'{a}}}1 {\'e}{{\'{e}}}1 {\'i}{{\'{i}}}1 {\'o}{{\'{o}}}1 {\'u}{{\'{u}}}1
    {\^a}{{\^{a}}}1 {\^e}{{\^{e}}}1 {\^o}{{\^{o}}}1
    {\~a}{{\~{a}}}1 {\~o}{{\~{o}}}1
    {\c c}{{\c{c}}}1
    {a\'}{{\'{a}}}1 {e\'}{{\'{e}}}1
    {i\'}{{\'{i}}}1 {o\'}{{\'{o}}}1 {u\'}{{\'{u}}}1
    {a\^}{{\^{a}}}1 {e\^}{{\^{e}}}1 {o\^}{{\^{o}}}1
    {a\~}{{\~{a}}}1 {o\~}{{\~{o}}}1
    {,}{{,}}1 {;}{{;}}1 {:}{{:}}1,
}

%% --- Bash style (without self-reference) -----------------------
\lstdefinestyle{bashstyle}{
  language=bash,
  basicstyle=\ttfamily\footnotesize,
  backgroundcolor=\color{codebg},
  frame=single,
  framerule=0.5pt,
  rulecolor=\color{codeframe},
  keywordstyle=\bfseries\color{pygray},
  commentstyle=\itshape\color{pygray},
  stringstyle=\color{pyorange},
  numbers=left,
  numberstyle=\tiny\color{pygray},
  numbersep=8pt,
  showstringspaces=false,
  breaklines=true,
  tabsize=2,
  captionpos=b,
  xleftmargin=1.5em,
  xrightmargin=0.5em,
  inputencoding=utf8,
  extendedchars=true,
  morekeywords={pip,pip3,python3,cd,jupyter,echo,source,export,conda,activate},
}

%% --- Highlight boxes --------------------------------------
\tcbuselibrary{skins, breakable}

\newtcolorbox{notebox}{
  colback=noteblue, colframe=noteborder,
  boxrule=0.8pt, arc=3pt,
  left=6pt, right=6pt, top=4pt, bottom=4pt,
  fonttitle=\bfseries, title=Note,
}

\newtcolorbox{warnbox}{
  colback=warnbg, colframe=warnborder,
  boxrule=0.8pt, arc=3pt,
  left=6pt, right=6pt, top=4pt, bottom=4pt,
  fonttitle=\bfseries, title={Warning},
}

%% --- Header / Footer --------------------------------------
\pagestyle{fancy}
\fancyhf{}
\fancyhead[L]{\small\texttt{econcomplex}}
\fancyhead[R]{\small Technical Documentation}
\fancyfoot[C]{\small\thepage}
\renewcommand{\headrulewidth}{0.4pt}

%% --- Sections ---------------------------------------------------
\titleformat{\section}{\large\bfseries\color{tblhead}}{}{0em}{}[\titlerule]
\titleformat{\subsection}{\normalsize\bfseries\color{pyblue}}{}{0em}{}
\titleformat{\subsubsection}{\normalsize\itshape}{}{0em}{}

%% --- Hyperref -------------------------------------------------
\hypersetup{
  colorlinks=true,
  linkcolor=pyblue,
  urlcolor=pyblue,
  citecolor=pygray,
  pdftitle={econcomplex -- Technical Documentation},
  pdfauthor={econcomplex contributors},
}

%% --- Macros ---------------------------------------------------
\newcommand{\pkg}[1]{\texttt{\textbf{#1}}}
\newcommand{\fn}[1]{\texttt{#1}}
\newcommand{\argm}[1]{\textit{\texttt{#1}}}
\newcommand{\mat}[1]{\mathbf{#1}}

%% ============================================================
\begin{document}

%% --- Cover page ----------------------------------------------------
\begin{titlepage}
  \centering
  \vspace*{3cm}
  {\Huge\bfseries\color{tblhead} econcomplex}\par
  \vspace{0.5cm}
  {\large\bfseries Technical Documentation --- Version 0.1.0}\par
  \vspace{1cm}
  \rule{10cm}{0.8pt}\par
  \vspace{0.8cm}
  {\large Python Library for Economic Complexity\\
  and Regional Science Indicators}\par
  \vspace{2cm}
  {\normalsize
    \textit{Elton Freitas}
  }\par
  \vspace{3cm}
  {\small\color{pygray}
    Dependencies: \texttt{numpy $\geq$ 1.21}, \texttt{pandas $\geq$ 1.3}, \texttt{scipy $\geq$ 1.7}\\[4pt]
    Compatible with Python 3.9+
  }\par
  \vfill
  {\small\color{pygray} March 2026}
\end{titlepage}

%% --- Table of contents --------------------------------------------------
\tableofcontents
\newpage

%% ============================================================
\section{Overview}

\pkg{econcomplex} is a Python library that consolidates the best
implementations of Economic Complexity and Regional Science indicators,
combining functionalities from four reference packages in the literature:

\begin{itemize}[leftmargin=1.5em]
  \item \textbf{EconGeo} (Balland, 2017) --- R package focused on geographic
        and regional specialization indicators
  \item \textbf{economiccomplexity} (Vargas Sep\'{u}lveda, 2020) --- R package
        with C++ backend (Armadillo) for complexity indicators
  \item \textbf{py-ecomplexity} (The Growth Lab at Harvard, 2020) --- complete
        Python pipeline with time series support
  \item \textbf{py-economic-complexity} --- modular Python implementation
        with Polars support and cross-space indicators
\end{itemize}

\subsection{Installation}

\begin{lstlisting}[style=bashstyle,
  caption={Installation via pip (local development)}]
# From the library root directory
pip install -e .

# Or adding the path manually
import sys
sys.path.insert(0, "/path/to/econcomplex")
import econcomplex as ec
\end{lstlisting}

\subsection{Module Structure}

\begin{center}
\begin{tabular}{lp{9cm}}
\toprule
\textbf{Module} & \textbf{Contents} \\
\midrule
\texttt{core}           & RCA, RPOP, Mcp, diversity, ubiquity, utilities \\
\texttt{complexity}     & ECI/PCI (3 methods), subnational ECI \\
\texttt{relatedness}    & Proximity, density, distance, co-occurrence, cross-space \\
\texttt{specialization} & LQ, Hachman, Krugman, specialization coef., similarity \\
\texttt{inequality}     & Gini, locational Gini, Hoover-Gini, HHI, Shannon entropy \\
\texttt{productivity}   & PRODY, EXPY, PGI, PEII \\
\texttt{patents}        & Ease of Recombination, Modular Complexity \\
\texttt{dynamics}       & Growth rates, entry/exit tracking \\
\texttt{outlook}        & COI, COG (Complexity Outlook) \\
\texttt{pipeline}       & \texttt{compute\_complexity()} function --- unified pipeline \\
\bottomrule
\end{tabular}
\end{center}

%% ============================================================
\section{Data Organisation}

This section describes the input format expected by the library,
ensuring that data is correctly structured before running any indicator.

% -- -- -- -- -- -- -- -- -- -- -- -- -- -- -- -- -- -- -- --
\subsection{Standard Format: Long Table (\textit{long format})}

The \pkg{econcomplex} library works exclusively with
\textbf{tables in long format} (\textit{tidy data}),
where each row represents \textbf{a unique combination} of
region $\times$ activity (and period, when applicable):

\begin{center}
\renewcommand{\arraystretch}{1.3}
\begin{tabular}{llll}
\toprule
\textbf{region} & \textbf{activity} & \textbf{value} & \textbf{year (optional)} \\
\midrule
\texttt{SP}  & \texttt{cnae\_10} & 12\,345 & 2022 \\
\texttt{SP}  & \texttt{cnae\_25} &  6\,789 & 2022 \\
\texttt{RJ}  & \texttt{cnae\_10} &  9\,012 & 2022 \\
\texttt{RJ}  & \texttt{cnae\_25} &  3\,456 & 2022 \\
\texttt{SP}  & \texttt{cnae\_10} & 13\,200 & 2023 \\
$\vdots$ & $\vdots$ & $\vdots$ & $\vdots$ \\
\bottomrule
\end{tabular}
\end{center}

\begin{warnbox}
  The library does \textbf{not} accept pivoted matrices (regions in rows,
  activities in columns). Use \fn{pivot\_to\_matrix()} only for
  low-level functions that explicitly require the matrix.
\end{warnbox}

% -- -- -- -- -- -- -- -- -- -- -- -- -- -- -- -- -- -- -- --
\subsection{Data Quality Rules}

\begin{enumerate}[leftmargin=1.5em]
  \item \textbf{No duplicate rows:} each (\textit{region},
        \textit{activity}) pair must appear exactly once per period.
  \item \textbf{Non-negative values:} the value column must contain
        only numbers $\geq 0$.
  \item \textbf{No \texttt{NaN}:} remove or impute missing values before
        passing the DataFrame.
  \item \textbf{Consistent granularity:} all regions must use
        the same geographic level (e.g., all municipalities
        \textit{or} all states).
  \item \textbf{Homogeneous activities:} use a single classification
        per analysis (e.g., CNAE 2-digit \textit{or} CNAE 4-digit,
        not mixed).
\end{enumerate}

% -- -- -- -- -- -- -- -- -- -- -- -- -- -- -- -- -- -- -- --
\subsection{Column Names --- \texttt{cols} Dictionary}

The pipeline uses a \argm{cols} dictionary to map the actual column
names of your DataFrame:

\begin{center}
\renewcommand{\arraystretch}{1.3}
\begin{tabular}{lll}
\toprule
\textbf{Key} & \textbf{Type} & \textbf{Description} \\
\midrule
\texttt{"loc"}  & required & Region / location column \\
\texttt{"act"}  & required & Activity / product / technology column \\
\texttt{"val"}  & required & Numeric value column (jobs, exports \ldots) \\
\texttt{"time"} & optional & Period column (year, quarter \ldots) \\
\bottomrule
\end{tabular}
\end{center}

\begin{lstlisting}[caption={Verifying and preparing the DataFrame before the pipeline}]
import pandas as pd
import econcomplex as ec

df = pd.read_csv("my_data.csv")

# Check structure
print(df.dtypes)
print(df.shape)
print(df.isnull().sum())        # check for NaN
print(df.duplicated().sum())    # check for duplicates

# Basic cleaning
df = df.dropna(subset=["region", "cnae", "jobs"])
df = df[df["jobs"] >= 0]        # remove negative values
df = df.drop_duplicates(subset=["region", "cnae", "year"])

# Run the pipeline
result = ec.compute_complexity(
    df,
    cols={"loc": "region", "act": "cnae",
          "val": "jobs", "time": "year"},
)
\end{lstlisting}

% -- -- -- -- -- -- -- -- -- -- -- -- -- -- -- -- -- -- -- --
\subsection{Supported File Formats}

\begin{center}
\renewcommand{\arraystretch}{1.3}
\begin{tabular}{lll}
\toprule
\textbf{Format} & \textbf{pandas reader} & \textbf{Notes} \\
\midrule
\texttt{.csv}     & \texttt{pd.read\_csv()}     & Most common; use \texttt{sep=";"} for Brazilian CSVs \\
\texttt{.parquet} & \texttt{pd.read\_parquet()} & Ideal for large datasets (RAIS, COMEX) \\
\texttt{.xlsx}    & \texttt{pd.read\_excel()}   & Requires \texttt{openpyxl}; avoid $>$100\,000 rows \\
\texttt{.feather} & \texttt{pd.read\_feather()} & High performance; requires \texttt{pyarrow} \\
\texttt{.dta}     & \texttt{pd.read\_stata()}   & Stata files; preserves labels \\
\bottomrule
\end{tabular}
\end{center}

\begin{lstlisting}[caption={Reading different file formats}]
import pandas as pd

# CSV with semicolon separator (Brazilian standard)
df = pd.read_csv("rais.csv", sep=";", encoding="latin-1",
                 usecols=["municipality", "cnae2d", "year", "jobs"],
                 dtype={"municipality": str, "cnae2d": str})

# Parquet (preferred for RAIS/COMEX)
df = pd.read_parquet("comex_ncm.parquet",
                     columns=["co_mun", "co_ncm", "year", "vl_fob"])

# Excel
df = pd.read_excel("data.xlsx", sheet_name="Sheet1",
                   usecols=["state", "sector", "year", "jobs"])
\end{lstlisting}

% -- -- -- -- -- -- -- -- -- -- -- -- -- -- -- -- -- -- -- --
\subsection{Examples of Compatible Datasets}

\begin{description}[leftmargin=1.5em, style=nextline]
  \item[\textbf{RAIS/MTE}]
    Formal employment by municipality $\times$ CNAE (2 or 4 digits)
    $\times$ year. Value column: active employment bonds on 31/12.
  \item[\textbf{COMEX/MDIC}]
    Exports by municipality $\times$ NCM $\times$ year.
    Value column: \texttt{VL\_FOB} (value in USD).
  \item[\textbf{INPI --- Patents}]
    Patent applications by State of the Federation $\times$
    IPC classification $\times$ year.
  \item[\textbf{Atlas of Economic Complexity (Harvard)}]
    Exports by country $\times$ product (HS) $\times$ year.
    Available at \url{https://atlas.cid.harvard.edu/}.
  \item[\textbf{Simulated data (testing)}]
    Use \fn{ec.make\_sample\_data(n\_locs=20, n\_acts=50, seed=42)}
    to generate a synthetic DataFrame.
\end{description}

%% ============================================================
\section{How to Run the Scripts}

This section describes all the ways to run the \pkg{econcomplex}
library, from direct execution of Python scripts to use in interactive
environments such as Jupyter and IDEs.

% -- -- -- -- -- -- -- -- -- -- -- -- -- -- -- -- -- -- -- --
\subsection{Prerequisites}

\begin{warnbox}
  Make sure the dependencies are installed before running any script.
\end{warnbox}

\begin{lstlisting}[style=bashstyle,
  caption={Installing the dependencies}]
# Install main dependencies
pip install numpy pandas scipy

# Install the library in development mode
pip install -e /path/to/econcomplex

# Verify installation
python3 -c "import econcomplex as ec; print(ec.__version__)"
\end{lstlisting}

% -- -- -- -- -- -- -- -- -- -- -- -- -- -- -- -- -- -- -- --
\subsection{Direct Execution via Terminal}

The simplest way to use the library is to run a \texttt{.py} script
directly in the terminal:

\begin{lstlisting}[style=bashstyle,
  caption={Running scripts in the terminal}]
# Navigate to the project folder
cd "/path/to/1 - Pos-doc research/data/script"

# Run the full example script
python3 econcomplex/example_usage.py

# Run your own script
python3 my_script.py
\end{lstlisting}

If the library is \textbf{not} installed via \texttt{pip},
add the path manually at the beginning of each script:

\begin{lstlisting}[caption={Adding path to sys.path (without pip install)}]
import sys
sys.path.insert(0, "/path/to/econcomplex")

import econcomplex as ec
print(ec.__version__)   # 0.1.0
\end{lstlisting}

% -- -- -- -- -- -- -- -- -- -- -- -- -- -- -- -- -- -- -- --
\subsection{Running in Jupyter Notebook / JupyterLab}

\begin{lstlisting}[style=bashstyle,
  caption={Opening Jupyter from the scripts folder}]
cd "/path/to/1 - Pos-doc research/data/script"
jupyter notebook
# or
jupyter lab
\end{lstlisting}

Inside the notebook, the first cell should configure the path:

\begin{lstlisting}[caption={Configuration cell in Jupyter}]
# Cell 1 -- configuration
import sys, os

BASE = os.path.expanduser(
    "~/Library/CloudStorage/"
    "OneDrive-Personnel/1 - Pos-doc research/data/script/econcomplex"
)
sys.path.insert(0, BASE)

import econcomplex as ec
import numpy as np
import pandas as pd

print("econcomplex", ec.__version__, "loaded successfully.")
\end{lstlisting}

\begin{lstlisting}[caption={Subsequent cells --- normal usage}]
# Cell 2 -- load data
df = pd.read_csv("my_data.csv")   # columns: region, sector, jobs

# Cell 3 -- compute complexity
result = ec.compute_complexity(
    df,
    cols={"loc": "region", "act": "sector", "val": "jobs"},
    method="eigenvector",
    compute_coi_cog=True,
)
result.head()

# Cell 4 -- display ECI by region
eci = (result[["region", "eci"]]
       .drop_duplicates()
       .set_index("region")
       .sort_values("eci", ascending=False))
print(eci)
\end{lstlisting}

% -- -- -- -- -- -- -- -- -- -- -- -- -- -- -- -- -- -- -- --
\subsection{Running in VS Code / PyCharm / Spyder}

In any IDE, configure the \textbf{Python interpreter} and the
\textbf{working directory} to point to the scripts folder.

\begin{lstlisting}[style=bashstyle,
  caption={Configure path in VS Code (settings.json)}]
// .vscode/settings.json
{
  "python.defaultInterpreterPath":
      "/Library/Frameworks/Python.framework/Versions/3.12/bin/python3",
  "terminal.integrated.cwd": "${workspaceFolder}"
}
\end{lstlisting}

In \textbf{Spyder}, go to \textit{Tools $\to$ PYTHONPATH Manager}
and add the \texttt{econcomplex/} directory.

% -- -- -- -- -- -- -- -- -- -- -- -- -- -- -- -- -- -- -- --
\subsection{Minimal Execution Script}

Ready-to-use template for any CSV-format dataset:

\begin{lstlisting}[caption={Minimal template --- copy and adapt}]
"""
Usage template for the econcomplex library.
Adapt the column names and the CSV file path.
"""
import sys
sys.path.insert(0, "/path/to/econcomplex")

import pandas as pd
import econcomplex as ec

# 1. Load data
df = pd.read_csv("data.csv")
print("Columns:", df.columns.tolist())
print("Shape:", df.shape)

# 2. Compute all indicators
result = ec.compute_complexity(
    df,
    cols={
        "loc":  "region_column_name",    # e.g.: "municipality", "state"
        "act":  "activity_column_name",  # e.g.: "cnae", "product"
        "val":  "value_column_name",     # e.g.: "jobs", "exports"
        "time": "year_column_name",      # optional
    },
    method="eigenvector",   # 'eigenvector' | 'reflections' | 'fitness'
    threshold=1.0,
    compute_coi_cog=True,
)

# 3. Export results
result.to_csv("complexity_results.csv", index=False)
print("Result saved to complexity_results.csv")
print("Generated columns:", sorted(result.columns.tolist()))

# 4. Individual indicators (optional)
mat = ec.pivot_to_matrix(df,
                         "region_column_name",
                         "activity_column_name",
                         "value_column_name")

eci, pci  = ec.eci_pci(mat)
hachman   = ec.hachman_index(mat)
krugman   = ec.krugman_index(mat)
entropy   = ec.shannon_entropy(mat)

output = pd.DataFrame({
    "eci": eci, "hachman": hachman,
    "krugman": krugman, "entropy": entropy,
})
print(output.round(3))
\end{lstlisting}

% -- -- -- -- -- -- -- -- -- -- -- -- -- -- -- -- -- -- -- --
\subsection{Running with RAIS / CNAE Data (real use case)}

\begin{lstlisting}[caption={Usage with RAIS data --- municipalities x CNAE}]
import sys
sys.path.insert(0, "/path/to/econcomplex")

import pandas as pd
import econcomplex as ec

# Load aggregated RAIS
# Expected columns: municipality, cnae2d, year, jobs
rais = pd.read_parquet("rais_municipality_cnae.parquet")

# Filter a single year for analysis
rais_2022 = rais[rais["year"] == 2022].copy()

# Pipeline
result = ec.compute_complexity(
    rais_2022,
    cols={"loc": "municipality", "act": "cnae2d", "val": "jobs"},
    method="eigenvector",
    compute_coi_cog=False,   # disable COG for large datasets
)

# Top 20 most complex municipalities
top20 = (result[["municipality", "eci"]]
         .drop_duplicates()
         .sort_values("eci", ascending=False)
         .head(20))
print(top20.to_string(index=False))

result.to_csv("complexity_municipalities_2022.csv", index=False)
\end{lstlisting}

% -- -- -- -- -- -- -- -- -- -- -- -- -- -- -- -- -- -- -- --
\subsection{Running a Multi-Year Panel}

To analyse the evolution of indicators over time,
simply include the time column in the \argm{cols} dictionary:

\begin{lstlisting}[caption={Panel execution --- multiple years}]
# rais with "year" column containing 2010, 2015, 2020
panel_result = ec.compute_complexity(
    rais,
    cols={
        "loc":  "municipality",
        "act":  "cnae2d",
        "val":  "jobs",
        "time": "year",       # activates the temporal loop
    },
    method="reflections",
    compute_coi_cog=False,
)

# The result already contains all years
print(panel_result["year"].value_counts().sort_index())

# Median ECI evolution by year
eci_trend = (panel_result[["year", "municipality", "eci"]]
             .drop_duplicates()
             .groupby("year")["eci"]
             .median()
             .round(4))
print(eci_trend)
\end{lstlisting}

\begin{notebox}
  The pipeline automatically iterates over each period, ensuring that
  the proximity matrix and complexity vectors are recalculated
  independently for each year.
\end{notebox}

%% ============================================================
\section{Core Indicators (\texttt{core})}

\subsection{Revealed Comparative Advantage (RCA)}

Balassa's RCA is the main specialization filter:

\[
  \text{RCA}_{il} = \frac{x_{il} / \sum_l x_{il}}
                         {\sum_i x_{il} / \sum_{il} x_{il}}
\]

\begin{lstlisting}[caption={Continuous and binary RCA}]
import econcomplex as ec
import pandas as pd

df = pd.read_csv("data.csv")
mat = ec.pivot_to_matrix(df, "region", "sector", "jobs")

rca      = ec.rca(mat)               # continuous values
mcp      = ec.mcp(mat, threshold=1)  # 1 if RCA >= 1, 0 otherwise
rca_rpop = ec.rpop(mat)              # Puga & Turrini RCA (continuous)
\end{lstlisting}

\subsection{Diversity and Ubiquity}

\begin{lstlisting}[caption={Diversity and ubiquity}]
div = ec.diversity(mat)    # number of sectors with RCA >= 1 per region
ubi = ec.ubiquity(mat)     # number of regions with RCA >= 1 per sector
\end{lstlisting}

%% ============================================================
\section{Economic Complexity (\texttt{complexity})}

\subsection{Eigenvector Method (Hidalgo \& Hausmann, 2009)}

\[
  \text{ECI} = \frac{\vec{v}_2 - \mu(\vec{v}_2)}{\sigma(\vec{v}_2)}, \quad
  \text{PCI} = \frac{\vec{u}_2 - \mu(\vec{u}_2)}{\sigma(\vec{u}_2)}
\]
where $\vec{v}_2$ is the second left eigenvector of matrix $\tilde{M}$:
\[
  \tilde{M}_{ij} = \frac{M_{ij}}{k_{c,0} \cdot k_{p,0}}
\]

\begin{lstlisting}[caption={ECI and PCI using the eigenvector method}]
eci, pci = ec.eci_pci(mat, method="eigenvector")
print(eci.head())   # Series indexed by region
print(pci.head())   # Series indexed by product/sector
\end{lstlisting}

\subsection{Method of Reflections (Hidalgo \& Hausmann, 2009)}

\begin{lstlisting}[caption={Method of Reflections}]
eci_r, pci_r = ec.eci_pci(mat, method="reflections", iterations=20)
\end{lstlisting}

\subsection{Fitness-Complexity (Tacchini et al., 2012)}

\[
  F_c^{(n+1)} = \sum_p M_{cp} Q_p^{(n)}, \quad
  Q_p^{(n+1)} = \frac{1}{\sum_c M_{cp} / F_c^{(n)}}
\]

\begin{lstlisting}[caption={Fitness-Complexity}]
fit, comp = ec.eci_pci(mat, method="fitness", iterations=100)
\end{lstlisting}

\subsection{Subnational ECI}

\begin{lstlisting}[caption={Subnational ECI with external PCI}]
# pci_national: PCI computed from national/international data
eci_sub = ec.subnational_eci(mat_sub, pci_national)
\end{lstlisting}

%% ============================================================
\section{Product Space and Relatedness (\texttt{relatedness})}

\subsection{Proximity ($\phi$)}

\[
  \phi_{ij} = \min\!\left(
    P(\text{RCA}_i \geq 1 \mid \text{RCA}_j \geq 1),\;
    P(\text{RCA}_j \geq 1 \mid \text{RCA}_i \geq 1)
  \right)
\]

\begin{lstlisting}[caption={Discrete and continuous proximity}]
phi_d = ec.proximity(mat, continuous=False)   # based on binary Mcp
phi_c = ec.proximity(mat, continuous=True)    # based on continuous RCA
\end{lstlisting}

\subsection{Density and Distance}

\[
  \omega_{ij} = \frac{\sum_k M_{kj} \phi_{ij}}{\sum_k \phi_{ij}}, \quad
  d_{ij} = 1 - \omega_{ij}
\]

\begin{lstlisting}[caption={Density and distance}]
density  = ec.density(mat, phi_d)   # DataFrame loc x act
distance = ec.distance(mat, phi_d)  # DataFrame loc x act
\end{lstlisting}

\subsection{Co-occurrence and Relatedness Indices}

\begin{lstlisting}[caption={Co-occurrence and relatedness indices}]
cooc    = ec.co_occurrence(mat)
phi_cos = ec.cosine_proximity(mat)     # cosine-based
phi_cor = ec.correlation_proximity(mat)
rca_rel = ec.relatedness(mat, phi_d)   # relatedness densities
\end{lstlisting}

\subsection{Cross-Space (relatedness between distinct classifications)}

\begin{lstlisting}[caption={Cross-space proximity and relatedness}]
# phi_cross: proximity between sectors and technologies (different matrices)
phi_cross = ec.cross_space_proximity(mat_sector, mat_tech)
density_c = ec.density(mat_sector, phi_cross)
\end{lstlisting}

%% ============================================================
\section{Regional Specialization (\texttt{specialization})}

\begin{lstlisting}[caption={Specialization indices}]
lq        = ec.location_quotient(mat)    # equivalent to RCA
hachman   = ec.hachman_index(mat)        # 0 (diversified) to 1 (specialized)
krugman   = ec.krugman_index(mat)        # structural difference between regions
coef_spec = ec.spec_coefficient(mat)     # specialization coefficient
sim_exp   = ec.export_similarity(mat)    # productive structure similarity
\end{lstlisting}

%% ============================================================
\section{Inequality and Concentration (\texttt{inequality})}

\begin{lstlisting}[caption={Inequality indicators}]
gini_r   = ec.gini(mat)                # Gini of value distribution
gini_loc = ec.locational_gini(mat)     # locational Gini (by sector)
hoover   = ec.hoover_gini(mat)         # Hoover index
hhi      = ec.hhi(mat)                 # Herfindahl-Hirschman
entropy  = ec.shannon_entropy(mat)     # Shannon entropy (by region)
\end{lstlisting}

%% ============================================================
\section{Productivity (\texttt{productivity})}

\begin{lstlisting}[caption={Productivity and product inequality}]
prody = ec.prody(mat, gdp_per_capita)  # income associated with the product
expy  = ec.expy(mat, prody)            # sophistication of the export basket
pgi   = ec.pgi(mat)                    # Product Gini Index
peii  = ec.peii(mat)                   # Product Export Inequality Index
\end{lstlisting}

%% ============================================================
\section{Patent-Based Complexity (\texttt{patents})}

\begin{lstlisting}[caption={Patent indicators (input: patent x technology matrix)}]
ease_rec = ec.ease_of_recombination(mat_pat)
mod_comp = ec.modular_complexity(mat_pat)
\end{lstlisting}

%% ============================================================
\section{Temporal Dynamics (\texttt{dynamics})}

\begin{lstlisting}[caption={Growth rates and entry/exit tracking}]
growth  = ec.growth_rates(df, "region", "sector", "jobs", "year")
entries = ec.entry_tracking(df, "region", "sector", "year")
exits   = ec.exit_tracking(df, "region", "sector", "year")
\end{lstlisting}

%% ============================================================
\section{Complexity Outlook (\texttt{outlook})}

The \textbf{COI} measures the latent diversification potential;
the \textbf{COG} measures the expected complexity gain:

\[
  \text{COI}_c = \sum_p (1 - M_{cp})\,\omega_{cp}\,Q_p
\]

\begin{lstlisting}[caption={COI and COG}]
coi = ec.coi(mat, phi_d, pci)
cog = ec.cog(mat, phi_d, pci)
\end{lstlisting}

%% ============================================================
\section{Unified Pipeline (\texttt{compute\_complexity})}

\subsection{Function Signature}

\begin{lstlisting}[caption={Pipeline function signature}]
def compute_complexity(
    df: pd.DataFrame,
    cols: dict,           # {"loc": str, "act": str, "val": str, "time": str}
    method: str = "eigenvector",   # 'eigenvector' | 'reflections' | 'fitness'
    threshold: float = 1.0,        # threshold for Mcp
    compute_coi_cog: bool = True,
    iterations: int = 20,          # for reflections/fitness
) -> pd.DataFrame:
    """
    Returns the original DataFrame enriched with columns:
    rca, mcp, diversity, ubiquity, eci, pci,
    density, distance, [coi, cog]
    """
\end{lstlisting}

\subsection{Parameters}

\begin{center}
\renewcommand{\arraystretch}{1.3}
\begin{tabular}{llp{7cm}}
\toprule
\textbf{Parameter} & \textbf{Type} & \textbf{Description} \\
\midrule
\texttt{df}              & \texttt{DataFrame} & Table in long format \\
\texttt{cols}            & \texttt{dict}      & Column mapping (see Section~2.3) \\
\texttt{method}          & \texttt{str}       & Complexity method \\
\texttt{threshold}       & \texttt{float}     & Mcp threshold (default: 1.0) \\
\texttt{compute\_coi\_cog} & \texttt{bool}    & Compute COI and COG \\
\texttt{iterations}      & \texttt{int}       & Iterations for iterative methods \\
\bottomrule
\end{tabular}
\end{center}

%% ============================================================
\section{Complete Examples}

\subsection{Example 1 --- Regional Employment Analysis}

\begin{lstlisting}[caption={Complete regional employment analysis}]
import sys
sys.path.insert(0, "/path/to/econcomplex")

import pandas as pd
import numpy as np
import econcomplex as ec

# Generate synthetic data (or load a real CSV)
df = ec.make_sample_data(n_locs=50, n_acts=30, seed=42)
# Columns: loc, act, val

# Full pipeline
result = ec.compute_complexity(
    df,
    cols={"loc": "loc", "act": "act", "val": "val"},
    method="eigenvector",
    compute_coi_cog=True,
)

# Summary of generated indicators
print(result.columns.tolist())
# ['loc','act','val','rca','mcp','diversity','ubiquity',
#  'eci','pci','density','distance','coi','cog']

# Top 10 regions by ECI
top_eci = (result[["loc","eci"]]
           .drop_duplicates()
           .sort_values("eci", ascending=False)
           .head(10))
print(top_eci)

# Save
result.to_csv("regional_result.csv", index=False)
\end{lstlisting}

\subsection{Example 2 --- Exports with COMEX Data}

\begin{lstlisting}[caption={Export analysis with COMEX/MDIC data}]
import sys
sys.path.insert(0, "/path/to/econcomplex")

import pandas as pd
import econcomplex as ec

# Load COMEX dataset aggregated by municipality x NCM x year
comex = pd.read_parquet("comex_mun_ncm.parquet")
comex_2023 = comex[comex["year"] == 2023].copy()

# Cleaning
comex_2023 = comex_2023.dropna(subset=["co_mun","co_ncm","vl_fob"])
comex_2023 = comex_2023[comex_2023["vl_fob"] > 0]

# Pipeline
result = ec.compute_complexity(
    comex_2023,
    cols={"loc": "co_mun", "act": "co_ncm", "val": "vl_fob"},
    method="eigenvector",
    compute_coi_cog=True,
)

# Save
result.to_csv("export_complexity_2023.csv", index=False)
print("Mean ECI:", result["eci"].dropna().mean().round(4))
\end{lstlisting}

\subsection{Example 3 --- RAIS Time Panel (2010--2022)}

\begin{lstlisting}[caption={Pipeline with panel data --- multiple periods}]
import sys
sys.path.insert(0, "/path/to/econcomplex")

import pandas as pd
import econcomplex as ec

rais = pd.read_parquet("rais_municipality_cnae_2010_2022.parquet")

# Temporal pipeline -- iterates over each year automatically
panel = ec.compute_complexity(
    rais,
    cols={
        "loc":  "municipality",
        "act":  "cnae2d",
        "val":  "jobs",
        "time": "year",
    },
    method="eigenvector",
    compute_coi_cog=False,
)

# ECI median evolution
eci_trend = (panel[["year","municipality","eci"]]
             .drop_duplicates()
             .groupby("year")["eci"]
             .median()
             .round(4))
print(eci_trend)

panel.to_csv("panel_complexity_2010_2022.csv", index=False)
\end{lstlisting}

%% ============================================================
\section{References}

\begin{description}[leftmargin=1.5em, style=nextline]

  \item[Balassa (1965)]
    Balassa, B. (1965). Trade liberalisation and ``revealed'' comparative
    advantage. \textit{The Manchester School}, 33(2), 99--123.

  \item[Hidalgo \& Hausmann (2009)]
    Hidalgo, C.A., \& Hausmann, R. (2009). The building blocks of economic
    complexity. \textit{PNAS}, 106(26), 10570--10575.

  \item[Balland (2017)]
    Balland, P.-A. (2017). \textit{EconGeo: Computing Key Indicators of the
    Geography of Innovation}. R package.

  \item[Vargas Sep\'{u}lveda (2020)]
    Vargas Sep\'{u}lveda, M. (2020). \textit{economiccomplexity: Computational
    Methods for Economic Complexity}. R package.

  \item[The Growth Lab at Harvard (2020)]
    The Growth Lab at Harvard University (2020). \textit{py-ecomplexity:
    Economic Complexity in Python}. GitHub.

  \item[Tacchini et al.\ (2012)]
    Tacchini, E., Loreto, V., Pietronero, L. (2012). A New Metrics for
    Countries' Fitness and Products' Complexity.
    \textit{Scientific Reports}, 2, 723.

\end{description}

%% ============================================================
\section{License and Citation}

This library is distributed under the MIT license. For use in academic
work, we suggest citing the original packages in the versions
corresponding to the indicators used.

\begin{lstlisting}[language={}, caption={Check the library version}]
import econcomplex as ec
print(ec.__version__)   # 0.1.0
\end{lstlisting}

\end{document}
